\chapter{Polynomials, factorization, and algebraic constructions}\label{chap:M07.11}

% Topic grouping: Formal coefficients and universal evaluation
\KEntry{K07.11.polynomial_definition}{Formal polynomials}
A polynomial over a unital ring $R$ is a sequence $(a_i)_{i\in\N}$
with finite support: only finitely many $a_i$ are nonzero. Its
coefficient of index $i$ is $a_i$. Equality of polynomials means equality
of every coefficient. The constant polynomial $C(a)$ has coefficient
$a$ at index zero and zero elsewhere; $X$ has coefficient $1$ at index
one and zero elsewhere. Thus
\[
 f=\sum_{i=0}^n C(a_i)X^i
\]
is a formal expression for its coefficient sequence, independent of
trailing zero coefficients. The map $C$ is injective because its
coefficient of index zero is its argument. A formal polynomial has not
yet been identified with any function on $R$.

\KEntry{K07.11.polynomial_arithmetic}{The polynomial ring}
Define addition coefficientwise, additive inverse coefficientwise, and
multiplication by convolution:
\[
 (f+g)_n=f_n+g_n,\qquad (fg)_n=\sum_{i+j=n}f_i g_j.
\]
These sequences still have finite support: supports of sums lie in
unions, and the support of a product lies in the set of pairwise sums
of support indices. The additive abelian-group laws follow coefficient
by coefficient. Both coefficients of $(fg)h=f(gh)$ are
$\sum_{i+j+k=n}f_i g_j h_k$, with the same order of factors, so
associativity in $R$ proves associativity here. Distributivity follows
by distributing each finite coefficient sum. The sequence $C(1)$ is
an identity, since only its coefficient of index zero contributes.
Consequently these operations make $R[X]$ a ring, and $C$ a unital
ring homomorphism. If $R$ is commutative, interchanging $i,j$ in the
product sum proves commutativity. Even if $R$ is not commutative,
$X C(a)=C(a)X$, since both have their sole possible nonzero
coefficient $a$ at index one; distributing gives $Xf=fX$ for all $f$.

\KEntry{K07.11.degree_conventions}{Degree and the zero polynomial}
For $f\ne0$ its degree is the largest index with nonzero coefficient,
and its leading coefficient $\operatorname{lc}(f)$ is that coefficient.
Set $\deg0=-\infty$ and $\operatorname{lc}(0)=0$; use
$-\infty+n=-\infty$. The natural degree agrees with the degree of a
nonzero polynomial, but is assigned the value zero at the zero
polynomial. Thus a statement about natural degrees of nonzero products
needs the stated nonvanishing hypotheses. Mathlib represents degree by
$\operatorname{WithBot}(\N)$ and calls the separate natural-valued
operation $\operatorname{natDegree}$.

A polynomial is monic if its leading coefficient is $1$. In a
nontrivial ring a monic polynomial is nonzero. In the zero ring the
zero polynomial is monic because $0=1$; this is why the division
theorems below explicitly require a nontrivial coefficient ring.
The constant polynomial $C(a)$ with $a\ne0$ has degree zero, including
when $a$ is a nonunit.

\KEntry{K07.11.degree_arithmetic}{Degree calculations}
\begin{proposition}\label{prop:polynomial-degree}
For polynomials over a commutative ring,
\[
 \deg(f+g)\le\max(\deg f,\deg g),\qquad
 \deg(fg)\le\deg f+\deg g.
\]
If $f,g\ne0$ and $\operatorname{lc}(f)\operatorname{lc}(g)\ne0$,
then $fg\ne0$, its degree is $\deg f+\deg g$, and its leading
coefficient is the product of the leading coefficients.
In particular a polynomial ring over a domain is a domain.
\end{proposition}
\begin{proof}
Above the maximum of the two support bounds both summands vanish.
If $m=\deg f$ and $n=\deg g$, every coefficient of $fg$ above $m+n$
vanishes, because in each summand either $i>m$ or $j>n$. At index
$m+n$, the only possibly nonzero summand has $i=m,j=n$ and is
$\operatorname{lc}(f)\operatorname{lc}(g)$. Its nonvanishing gives
the equality assertion. If a factor is zero, the inequalities follow
from the convention for $-\infty$. In a domain two nonzero leading
coefficients have nonzero product, so two nonzero polynomials cannot
have zero product.
\end{proof}
The condition on leading coefficients also holds if one is a unit.
With zero divisors the product degree can drop: in $(\Z/4\Z)[X]$,
$(2X)^2=0$, although both factors have degree one.

\KEntry{K07.11.evaluation_universal}{The universal property of evaluation}
\begin{theorem}\label{thm:polynomial-evaluation}
Let $\varphi:R\to S$ be a unital ring homomorphism, and let $s\in S$
commute with every $\varphi(a)$, $a\in R$. There is exactly one
unital ring homomorphism $\Phi:R[X]\to S$ such that
$\Phi(C(a))=\varphi(a)$ and $\Phi(X)=s$. It is
\[
 \Phi\left(\sum_i C(a_i)X^i\right)=\sum_i\varphi(a_i)s^i.
\]
\end{theorem}
\begin{proof}
The finite sum defines a function of the coefficient sequence and
preserves addition and one. Its value on a product is
$\sum_{i,j}\varphi(a_i)\varphi(b_j)s^{i+j}$.
Since $s^i$ commutes with $\varphi(b_j)$, this is
$\sum_{i,j}(\varphi(a_i)s^i)(\varphi(b_j)s^j)$, the product
of the two evaluations. The required images of constants and $X$
follow directly. Any homomorphism with those images must send the
finite expression for a polynomial to the displayed sum, proving
uniqueness.
\end{proof}
For commutative $S$ the commuting condition is automatic. It cannot be
discarded for a noncommutative target, because $XC(a)=C(a)X$ would
force it. Ordinary evaluation $f(a)$ uses $\varphi=\id_R$ in a
commutative ring. Substitution $f(g)$ uses target $R[X]$.

\KEntry{K07.11.coefficient_map}{Changing the coefficient ring}
A unital homomorphism $\varphi:R\to S$ induces a unital homomorphism
$\varphi_*:R[X]\to S[X]$ by applying $\varphi$ to every coefficient.
Indeed convolution and finite sums commute with $\varphi$; alternatively
apply Theorem~\ref{thm:polynomial-evaluation} to $C\varphi$ and $X$.
It sends $C(a)$ to $C(\varphi(a))$ and fixes $X$.
If $\varphi$ is injective, a coefficient is zero precisely when its
image is zero, so $\varphi_*$ is injective and preserves degree and
leading coefficients. A noninjective map may kill a leading
coefficient: $2X+1$ becomes $1$ on reducing modulo $2$.

For an ideal $I$ of a commutative ring let
$I[X]=\{f:\text{every coefficient of }f\text{ belongs to }I\}$.
The coefficient formulas show that it is an ideal. The map
$R[X]\to(R/I)[X]$ is onto: lift the finitely many coefficients of
each target polynomial and set all others to zero. Its kernel is
exactly $I[X]$. The first isomorphism theorem therefore gives
\begin{equation}\label{eq:polynomial-coefficient-quotient}
 R[X]/I[X]\ \cong\ (R/I)[X].
\end{equation}
In particular $I[X]$ is also the ideal generated by the constants
$C(a)$ for $a\in I$: one inclusion follows from the ideal property,
and the other follows by writing $f=\sum C(a_i)X^i$.

\KEntry{K07.11.polynomial_units}{Units and nilpotent coefficients}
\begin{proposition}\label{prop:polynomial-units}
Over a commutative domain, the units of $R[X]$ are precisely the
constant polynomials whose coefficients are units of $R$.
Over any commutative ring, $f=\sum a_iX^i$ is a unit if and only if
$a_0$ is a unit and every $a_i$ with $i>0$ is nilpotent.
\end{proposition}
\begin{proof}
In a domain $fg=1$ gives $\deg f+\deg g=0$, so both degrees are zero;
the constant coefficients multiply to one. Conversely a unit constant
has its constant inverse.

In general, if $fg=1$, evaluation at zero gives $a_0g(0)=1$.
For each prime ideal $P$, the image of $f$ in $(R/P)[X]$ is a unit
over a domain, so all its positive-degree coefficients vanish.
Thus each $a_i$, $i>0$, belongs to every prime ideal, and
Lemma~\ref{lem:nilpotent-prime-separation} makes it nilpotent.
This implication also holds in the zero ring, whose elements are
all nilpotent.

Conversely suppose $a_0$ is a unit and choose exponents
$e_i\ge1$ with $a_i^{e_i}=0$ for the finitely many positive
coefficients. The ideal $J$ they generate satisfies $J^N=0$ for
$N=1+\sum_i(e_i-1)$: expanding a product of $N$ generators,
some generator occurs at least $e_i$ times. Write
$f=C(a_0)(1+h)$, where all coefficients of $h$ belong to $J$.
The coefficients of $h^N$ are sums of products from $J^N$, so
$h^N=0$. Consequently
$(1+h)\sum_{j=0}^{N-1}(-h)^j=1$ gives a polynomial inverse for $f$.
If there are no positive coefficients, use $N=1$ and $h=0$.
\end{proof}
For example $1+2X$ is a nonconstant unit in $(\Z/4\Z)[X]$,
with inverse $1-2X$. The domain hypothesis in the first assertion
is therefore substantial.

% Topic grouping: Division, roots, and multiplicities
\KEntry{K07.11.monic_division}{Division by a monic polynomial}
\begin{theorem}\label{thm:monic-polynomial-division}
Let $R$ be a nontrivial commutative ring and $g\in R[X]$ monic.
Every $f$ has unique $q,r$ such that
\[
 f=qg+r,\qquad \deg r<\deg g.
\]
The same conclusion holds when the leading coefficient of $g$ is a
unit.
\end{theorem}
\begin{proof}
First assume $g$ monic, of degree $m$. If $f=0$ or $\deg f<m$,
take $q=0,r=f$. Otherwise write $n=\deg f$ and $a=\operatorname{lc}(f)$.
The polynomial $f-C(a)X^{n-m}g$ has degree less than $n$: terms
above $n$ are zero and the coefficient of index $n$ cancels.
Inducting on the nonnegative degree produces its quotient and
remainder; adding $C(a)X^{n-m}$ to the quotient gives the required
identity for $f$. This induction terminates because a strictly
decreasing sequence of natural numbers terminates.

If $qg+r=q'g+r'$ and $q-q'\ne0$, the monic leading coefficient
gives
$\deg((q-q')g)=\deg(q-q')+m\ge m$ by
Proposition~\ref{prop:polynomial-degree}. But
$(q-q')g=r'-r$ has degree less than $m$, a contradiction.
Thus $q=q'$ and then $r=r'$.
For unit leading coefficient $u$, divide by the monic polynomial
$u^{-1}g$ and absorb $u^{-1}$ into the quotient.
\end{proof}
When $g=1$ the remainder must be zero. Over $\Z$, division of $X$
by the constant polynomial $2$ with remainder degree less than zero
would require $X=2q$, impossible at its coefficient of index one.
Thus nonzero divisors alone do not suffice over a coefficient ring.

\KEntry{K07.11.field_division}{Division over a field}
\begin{corollary}\label{thm:polynomial-division}
For a field $F$, a nonzero $g\in F[X]$, and any $f\in F[X]$,
there are unique $q,r$ with $f=qg+r$ and
$r=0$ or $\deg r<\deg g$.
\end{corollary}
\begin{proof}
The nonzero leading coefficient of $g$ is a unit, so apply the preceding
theorem. Its degree condition includes $r=0$ through $\deg0=-\infty$.
\end{proof}
This gives an algorithm using only field operations: cancel the
highest term until the remainder degree is small. A nonzero constant
divisor has remainder zero; no assertion of division by zero is made.

\KEntry{K07.11.field_polynomial_pid}{Polynomial rings over fields}
\begin{theorem}\label{thm:field-polynomial-pid}
For a field $F$, $F[X]$ is Euclidean with degree as the decreasing
quantity for nonzero remainders. Consequently it is a PID and a UFD.
Conversely, if $R$ is a commutative ring and $R[X]$ is a PID in the
domain sense, then $R$ is a field.
\end{theorem}
\begin{proof}
The polynomial ring is a domain by Proposition~\ref{prop:polynomial-degree},
and division with remainder supplies the Euclidean condition.
Theorems~\ref{thm:euclidean-implies-pid} and~\ref{thm:pid-implies-ufd}
give the two consequences.

For the converse, the injective constants map embeds $R$ in the
domain $R[X]$, so $R$ is a domain. Evaluation at zero is onto with
kernel $(X)$: a polynomial has zero constant coefficient precisely
when it is $X$ times the polynomial formed by shifting its other
coefficients down by one. Hence $R[X]/(X)\cong R$ is a domain,
so $(X)$ is prime. It is a nonzero proper ideal, since $X\ne0$
and evaluation at zero sends it to zero but sends $1$ to $1$.
In a PID every nonzero prime ideal is maximal by
Proposition~\ref{prop:pid-primes}. Thus its quotient $R$ is a field
by Proposition~\ref{prop:maximal-quotient}.
\end{proof}

\KEntry{K07.11.factor_remainder}{The factor and remainder theorems}
\begin{proposition}\label{prop:factor-remainder}
Let $R$ be commutative and $a\in R$. The remainder of $f$ on division
by $X-C(a)$ is $C(f(a))$, and
$f(a)=0$ if and only if $X-C(a)$ divides $f$.
If $R$ is a domain, $X-C(a)$ is prime and irreducible.
\end{proposition}
\begin{proof}
In a nontrivial ring the divisor is monic of degree one, so the remainder
is a constant $C(b)$. Evaluation of
$f=q(X-C(a))+C(b)$ at $a$ gives $b=f(a)$. The factor equivalence
follows, in one direction by this formula and in the other by evaluating
a product. In the zero ring every assertion of equality and divisibility
here holds directly; no division theorem is needed.

For a domain evaluation $R[X]\to R$ is onto and has kernel
$(X-C(a))$ by the factor equivalence. The quotient is the domain $R$,
so the generated ideal is prime. The generator is nonzero and nonunit
by its degree, and the prime-ideal test makes it a prime element.
Proposition~\ref{prop:prime-implies-irreducible} proves irreducibility.
\end{proof}

\KEntry{K07.11.root_bound}{The number of roots}
\begin{theorem}\label{thm:polynomial-root-bound}
A nonzero polynomial of degree $n$ over a domain has at most $n$
distinct roots. More strongly the sum of its root multiplicities is
at most $n$.
\end{theorem}
\begin{proof}
If $a$ is a root, write $f=(X-a)g$ by the factor theorem.
Then $g\ne0$ and $\deg g=n-1$ by the degree formula.
If $b\ne a$ is another root, $0=f(b)=(b-a)g(b)$ and cancellation
gives $g(b)=0$. Thus repeatedly removing distinct roots decreases
degree each time and cannot remove more than $n$ factors.
This argument also proves finiteness: if more than $n$ distinct roots
existed, choose $n+1$ of them and obtain a negative degree for a
nonzero quotient.

For multiplicities, remove $X-a$ as often as possible at each root;
each removal lowers degree by one, so the process stops. Removing a
factor at $a$ leaves the multiplicity at a different $b$ unchanged.
Indeed $X-b$ is prime by Proposition~\ref{prop:factor-remainder} and
does not divide $X-a$, since its value at $b$ is $b-a\ne0$.
Repeated prime cancellation therefore says
$(X-b)^k\mid(X-a)g$ if and only if $(X-b)^k\mid g$.
The total number of removed linear factors is consequently the sum
of the original multiplicities, and degree bounds it by $n$.
\end{proof}
The zero polynomial is excluded: it vanishes at every element, and
every power of every linear polynomial divides it. A library convention
may give the root multiset of zero or a numerical zero multiplicity a
default value; that convention is not this finite-root theorem.

\KEntry{K07.11.multiplicity_derivative}{Formal derivatives and repeated roots}
For a nonzero polynomial over a field the multiplicity of $a$ is the
largest $m\ge0$ for which $(X-a)^m\mid f$. This maximum exists by
the degree bound on such divisors. Thus
$f=(X-a)^m g$ with $g(a)\ne0$: otherwise the factor theorem would
increase $m$.

Define the formal derivative by
$f'=\sum_{i\ge1} i a_i X^{i-1}$, where the integer $i$ acts by
repeated addition. It is additive. For monomials, the identity
$(X^{i+j})'=(i+j)X^{i+j-1}$ equals the sum from differentiating
the two factors. Distributing a finite product of sums proves
$(fg)'=f'g+fg'$ for arbitrary polynomials.
Induction gives
\[
 f'=m(X-a)^{m-1}g+(X-a)^m g'\qquad(m\ge1).
\]
If $m=1$, evaluation at $a$ is $g(a)\ne0$; if $m\ge2$ it is zero.
Therefore a root has multiplicity greater than one exactly when
$f(a)=f'(a)=0$. If $f(a)\ne0$, the multiplicity is zero, so the
same pair of equalities characterizes multiplicity at least two.
This uses no analytic limiting operation. In characteristic $p$,
$(X^p)'=pX^{p-1}=0$ although $X^p$ is nonconstant. The zero
polynomial has no largest multiplicity and is excluded.

\KEntry{K07.11.low_degree_irreducible}{Small degrees}
Every degree-one polynomial over a field is irreducible: a factorization
has degrees summing to one, so one factor has degree zero and is a unit.
For degree two or three, a reducible polynomial is a product of two
positive-degree polynomials, and their degrees force one to have degree
one. Such a factor is a nonzero scalar multiple of $X-a$ for some
field element $a$, giving a root. Conversely a root gives a linear
factor and a quotient of positive degree, so gives reducibility.
Hence in these two degrees irreducibility is equivalent to having no
root. In higher degrees this test fails: $(X^2+1)^2$ has no rational
root, since a rational square is nonnegative, but it is reducible in
$\Q[X]$.

\KEntry{K07.11.rational_root}{The rational root test}
Let $0\ne f=a_nX^n+\cdots+a_0\in\Z[X]$, $a_n\ne0$.
If $r/s$ is a root, with $\gcd(r,s)=1$ and $s\ne0$, then
$r\mid a_0$ and $s\mid a_n$.
There is no root if $n=0$. For $n\ge1$, multiply the root equation
by $s^n$:
\[
 a_nr^n+a_{n-1}r^{n-1}s+\cdots+a_0s^n=0.
\]
It follows that $r\mid a_0s^n$ and $s\mid a_nr^n$.
Coprimality gives $\gcd(r,s^n)=\gcd(s,r^n)=1$:
a prime dividing one of those gcds would divide both $r$ and $s$,
using prime divisibility of a power. Bezout's identity then cancels
the respective coprime factor, proving both divisibilities.
If $r=0$, reduced form forces $s=\pm1$, the equation says $a_0=0$,
and the same conclusions hold, including $0\mid0$.
For a monic polynomial $s\mid1$, so every rational root is an integer.

% Topic grouping: Content, Gauss's lemma, and irreducibility tests
\KEntry{K07.11.content_primitive}{Content and primitive polynomials}
Assume $R$ is a UFD. The content $c(f)$ is a gcd of the finitely many
coefficients of $f$, determined up to multiplication by a unit.
Set $c(0)=0$. A nonzero polynomial is primitive if its coefficients
have unit gcd, or equivalently no prime element divides all its
coefficients. Indeed a nonunit common divisor has a prime factor,
and a common prime is itself a nonunit common divisor.
For $f\ne0$, divide all coefficients by $c(f)$ to obtain
\[
 f=C(c(f))f_0.
\]
This $f_0$ is primitive. To see this precisely, for a prime $p$ write
$v_p(a)$ for its exponent in a nonzero element and $v_p(0)=+\infty$.
The exponent of $p$ in the gcd of a finite collection not all zero is
the minimum of its coefficient exponents. Division by their gcd
subtracts that minimum, leaving minimum zero for every prime.
This also constructs the gcd from unique factorization.
The primitive part $f_0$ depends on the chosen associate of the content.
No primitive part of zero is needed in these arguments; conventions
that assign a value there do not make zero primitive.
Every monic polynomial is primitive because a divisor of all
coefficients must divide its leading coefficient $1$.

\KEntry{K07.11.gauss_content}{Products and content}
\begin{theorem}\label{thm:gauss-content}
Over a UFD the product of two primitive polynomials is primitive.
For arbitrary polynomials,
\[
 c(fg)\ \text{is associated to}\ c(f)c(g).
\]
\end{theorem}
\begin{proof}
Suppose primitive $f,g$ had nonprimitive product. A prime $p$ would
divide all coefficients of $fg$. In $(R/(p))[X]$ their images would
have zero product. The coefficient ring $R/(p)$ is a domain, since
$p$ is prime; its polynomial ring is a domain. Therefore one image
is zero, meaning $p$ divides all coefficients of one factor,
contrary to primitivity.

For nonzero $f,g$ write $f=C(a)f_0$, $g=C(b)g_0$ with
$a=c(f),b=c(g)$ and primitive $f_0,g_0$. Their product $h=f_0g_0$
is primitive. For each prime $p$ the minimum exponent among the
coefficients of $C(ab)h$ equals
$v_p(a)+v_p(b)+0$. Thus its content has exactly the prime
exponents of $ab$, proving association by unique factorization.
If one factor is zero, both compared contents are zero and the
assertion holds directly.
\end{proof}
The invariant assertion is association. Literal equality requires a
chosen multiplicative normalization of associates, for instance
nonnegative integer gcds over $\Z$ or monic gcds over a field;
an arbitrary selection of associates need not be multiplicative.

\KEntry{K07.11.gauss_fraction_irreducible}{Gauss's lemma over the fraction field}
\begin{lemma}\label{lem:primitive-scalar}
Let $K=\Frac(R)$ for a UFD $R$. Every nonzero $A\in K[X]$ has
the form $uA_0$, where $u\in K^\times$ and $A_0\in R[X]$ is
primitive. If $h\in R[X]$ is primitive and $t\in K^\times$
satisfies $th\in R[X]$, then $t\in R$. If $th$ is primitive
as well, $t$ is a unit of $R$.
\end{lemma}
\begin{proof}
There are only finitely many coefficients of $A$. Multiplying their
nonzero denominators gives $d\ne0$ with $dA\in R[X]$.
Extracting its content gives $dA=cA_0$, hence $A=(c/d)A_0$.

Write $t=r/s$ with $r,s\ne0$ and with their common prime factors
cancelled, so $\gcd(r,s)$ is a unit. Integrality of each coefficient
$r h_i/s$ says $s\mid rh_i$. If a prime $p$ divided $s$, its
full exponent in $s$ would have to occur in $h_i$, since it does
not occur in $r$. In particular $p$ would divide every $h_i$,
contrary to primitivity. Thus $s$ is a unit and $t\in R$.
The content of $th$ is then associated to $t$, by the coefficient
exponent calculation above. If $th$ is primitive, this makes $t$
a unit.
\end{proof}

\begin{theorem}\label{thm:gauss-fraction}
A primitive polynomial in $R[X]$, for $R$ a UFD, is irreducible
over $R$ if and only if its image in $\Frac(R)[X]$ is irreducible.
More generally, any factorization of a nonzero $f\in R[X]$ into two
positive-degree factors over the fraction field can be rescaled into
such a factorization over $R$.
\end{theorem}
\begin{proof}
For a factorization $f=AB$ over $K$, write $A=uA_0$, $B=vB_0$
as in the lemma. The product $A_0B_0$ is primitive, and
$f=(uv)A_0B_0$ is integral. The lemma gives $uv\in R$, so
$f=((uv)A_0)B_0$ is the required integral factorization.
Nonzero scalar rescaling preserves both degrees.
If $f$ is primitive, the same lemma says $uv$ is a unit of $R$.

For a primitive polynomial of positive degree, a nonunit constant
factor over $R$ would divide all coefficients, which is impossible.
Thus reducibility over $R$ is equivalent to a factorization into
two positive-degree factors. The same equivalence holds over the
field $K$, by the description of its polynomial units.
The rescaling result proves one direction; the coefficient embedding
and degree preservation prove the other. A primitive constant is
a unit in both rings, so is irreducible in neither. These cases
exhaust primitive polynomials.
\end{proof}
In particular the theorem applies to monic polynomials. Without
primitivity it fails: $2X+2$ is reducible over $\Z$ because of its
nonunit content, but is a degree-one irreducible over $\Q$.

\KEntry{K07.11.polynomial_ufd}{Unique factorization is inherited}
\begin{theorem}\label{thm:polynomial-ufd}
For a commutative domain $R$, the ring $R$ is a UFD if and only if
$R[X]$ is a UFD. Hence a polynomial ring in finitely many variables
over a UFD is a UFD.
\end{theorem}
\begin{proof}
First suppose $R$ is a UFD and put $K=\Frac(R)$.
Factor a nonzero $f\in R[X]$ as its content times a primitive
polynomial $f_0$. The content factors into primes in $R$.
The field polynomial ring $K[X]$ is a UFD by
Theorem~\ref{thm:field-polynomial-pid}. Factor $f_0$ there, and
rescale each positive-degree irreducible factor to a primitive
polynomial in $R[X]$ using Lemma~\ref{lem:primitive-scalar}.
Their product is primitive, so its scalar multiplier in $f_0$ is a
unit of $R$, again by that lemma. Each rescaled factor is irreducible
in $R[X]$ by Theorem~\ref{thm:gauss-fraction}. If $f_0$ is constant,
it is already a unit. We have proved factorization existence in all
cases.

We next show all factors just used are prime in $R[X]$, which will
give uniqueness. A constant prime $p$ remains prime, since
$R[X]/(C(p))\cong(R/(p))[X]$ by
\eqref{eq:polynomial-coefficient-quotient}; here $(C(p))=(p)[X]$
by the finite coefficient expression. The quotient is a domain.

Let $h\in R[X]$ be primitive and irreducible of positive degree.
It is prime in $K[X]$, since that ring is a PID. If $h\mid ab$
in $R[X]$, it therefore divides $a$ or $b$ in $K[X]$.
Suppose $a=hQ$. If $a=0$, divisibility in $R[X]$ is immediate.
Otherwise write $Q=tQ_0$ with $Q_0$ primitive. The product
$hQ_0$ is primitive and $a=t(hQ_0)$ is integral, so the scalar
lemma gives $t\in R$. Thus $a=h(tQ_0)$ in $R[X]$.
The case for $b$ is identical. Hence $h$ is prime in $R[X]$.
All irreducible factors in our existence factorization are therefore
prime. Cancellation of a prime from two factorizations, followed by
induction, proves uniqueness up to units and permutation exactly as
in Lemma~\ref{lem:prime-factorization-unique}.

Conversely suppose $R[X]$ is a UFD. A nonzero constant $a$ factors
there into irreducibles. Degree additivity in a domain forces all
factors to be constants. A constant is a polynomial unit exactly
when its coefficient is a unit of $R$; and a factorization of a
nonzero constant in $R[X]$ must have constant factors. Consequently
a constant is irreducible in $R[X]$ exactly when its coefficient
is irreducible in $R$. Existence and uniqueness for the polynomial
factorization of $C(a)$ give existence and uniqueness in $R$.
Finally iterate the one-variable result to obtain finitely many
variables, starting with the zero-variable ring $R$.
\end{proof}

\KEntry{K07.11.reduction_irreducibility}{Reduction modulo a prime}
\begin{proposition}\label{prop:irreducible-reduction}
Let $\varphi:R\to S$ be a unital homomorphism between domains.
If $f\in R[X]$ is monic and $\varphi_*(f)$ is irreducible in
$S[X]$, then $f$ is irreducible in $R[X]$.
In particular a monic integer polynomial with irreducible reduction
modulo a prime $p$ is irreducible over both $\Z$ and $\Q$.
\end{proposition}
\begin{proof}
The image is a nonunit of positive degree, so the same holds for the
monic $f$. If $f=gh$ with both factors nonunits, their nonzero
leading coefficients multiply to $1$. They are therefore units.
A factor of degree zero would then be a unit, so both degrees are
positive. If $u=\operatorname{lc}(g)$, replace the factors by
$u^{-1}g$ and $uh$ to make both monic. Their images under
$\varphi_*$ remain monic of the same positive degrees, contradicting
irreducibility of $\varphi_*(f)$.
For reduction modulo $p$, take $S=\F_p$ and use Gauss's lemma for
the monic, hence primitive, integer polynomial.
\end{proof}
Outside the monic case both primitivity and degree preservation matter:
$2X+2$ reduces to the degree-one irreducible $2X+2$ modulo $3$,
but is reducible over $\Z$ because it is not primitive. For a
primitive nonmonic polynomial a reduction test is valid if its
leading coefficient survives; a positive-degree factorization then
retains the degrees of both factors, since their leading-coefficient
product survives in a domain. The converse test fails: $X^2+1$ is
irreducible over $\Q$ but reduces to $(X+1)^2$ modulo $2$.

\KEntry{K07.11.eisenstein}{Eisenstein's criterion}
\begin{theorem}\label{thm:eisenstein}
Let $R$ be a domain, $P$ a prime ideal, and $f\in R[X]$ monic
of positive degree $n$. Suppose every nonleading coefficient belongs
to $P$ and the constant coefficient does not belong to $P^2$.
Then $f$ is irreducible in $R[X]$.
If $R$ is a UFD, the same holds for a primitive nonmonic $f$
whose leading coefficient is outside $P$. In either case, when
$R$ is a UFD, $f$ is also irreducible over $\Frac(R)$.
\end{theorem}
\begin{proof}
Suppose first $f$ is monic and has a factorization into nonunits.
As in the reduction proposition, the factors may be made monic
and have positive degrees, say $f=gh$. Reduce coefficients modulo
$P$ to the domain $D=R/P$. We obtain
$\bar g\bar h=X^n$, with degrees preserved. Embed $D$ in its
fraction field $L$. In the UFD $L[X]$, the polynomial $X$ is
prime and the only irreducible factor of $X^n$. Thus the monic
positive-degree factors are $\bar g=X^i$, $\bar h=X^j$ with
$i,j>0$. To justify this use of uniqueness, factor either factor
in $L[X]$ and compare its irreducible factors with the product of
$n$ copies of $X$; its scalar factor is $1$ by monicity.
Both constant coefficients are consequently zero in $D$, so the
constant coefficients of $g,h$ belong to $P$. Their product,
the constant coefficient of $f$, belongs to $P^2$, a contradiction.

In the nonmonic case primitivity rules out any nonunit constant factor.
The leading coefficient outside $P$ means that in a proper
positive-degree factorization both leading coefficients are outside
$P$, since $P$ is prime. Degrees again survive reduction, and over
$L$ the product is a nonzero scalar times $X^n$. Unique
factorization makes each positive-degree factor a scalar times a
positive power of $X$. The identical constant-coefficient argument
gives the contradiction. Finally Gauss's lemma applies to the
primitive polynomial, including the monic case.
\end{proof}
For $R=\Z$ and $P=(p)$ this says that $p$ divides all nonleading
coefficients, $p$ does not divide the leading coefficient, and $p^2$
does not divide the constant coefficient, together with primitivity
if the leading coefficient is not $1$. For instance $X^n-p$ is
irreducible over $\Q$ for every $n\ge1$.

\KEntry{K07.11.shift_irreducibility}{A change of variable}
For a commutative ring, $a\in R$, and $u\in R^\times$, substitution
$X\mapsto uX+a$ defines a ring automorphism fixing constants.
Its inverse substitutes $u^{-1}(X-a)$: either composite fixes
constants and sends $X$ to $X$, and is therefore the identity by
Theorem~\ref{thm:polynomial-evaluation}. An automorphism preserves
units and all factorizations, so preserves irreducibility in both
directions. In particular shifts $X\mapsto X+a$ preserve it.
For example
\[
 (X+1)^4+1=X^4+4X^3+6X^2+4X+2
\]
is Eisenstein at $2$, since its constant coefficient is not divisible
by $4$. It is irreducible in $\Q[X]$, so the inverse shift proves
$X^4+1$ irreducible there. A successful shift can make a criterion
applicable without changing the factorization question.

% Topic grouping: Quotients, functions, and the boundary of the chapter
\KEntry{K07.11.irreducible_quotient}{Adjoining a root by a quotient}
\begin{theorem}\label{thm:irreducible-polynomial-quotient}
For a field $F$ and a nonconstant $f\in F[X]$, the quotient
$F[X]/(f)$ is a field if and only if $f$ is irreducible.
In this quotient the class $\alpha$ of $X$ satisfies $f(\alpha)=0$,
where the coefficients are mapped to their constant classes.
\end{theorem}
\begin{proof}
In the PID $F[X]$, a nonzero nonunit $f$ is irreducible exactly
when $(f)$ is maximal by Proposition~\ref{prop:pid-primes}.
The maximal-quotient criterion proves the equivalence.
The quotient map agrees with evaluation at $\alpha$ on constants
and on $X$, hence everywhere by the universal property; its value
at $f$ is zero because $f\in(f)$.

One can see the obstruction for a reducible $f$ explicitly.
Write $f=gh$ with both factors of positive degree less than $\deg f$.
Neither class is zero: if $g=af$ with $a\ne0$, degree additivity
would give $\deg g\ge\deg f$, and $a=0$ would give $g=0$.
The same applies to $h$. Their classes are nonzero and have zero
product, so the quotient has zero divisors.
\end{proof}
The restrictions distinguish two boundary cases. If $f=0$, the
quotient is $F[X]$, which is a domain but not a field because $X$
is a nonunit. If $f$ is a nonzero constant, it is a unit and its
quotient is the zero ring, which is not a field under our convention.
For example $\R[X]/(X^2+1)$ is a field: its generator $\alpha$
satisfies $\alpha^2=-1$, and the map $a+b\alpha\mapsto a+bi$
identifies it with $\C$. The normal forms and injectivity in this
example follow from the next theorem, or directly from division
by $X^2+1$ and equality of real and imaginary parts.

\KEntry{K07.11.quotient_normal_form}{Normal forms and dimension}
\begin{theorem}\label{thm:polynomial-quotient-basis}
Let $0\ne f\in F[X]$ have degree $n$. Every class in $F[X]/(f)$
has a unique representative of degree less than $n$.
The classes of $1,X,\ldots,X^{n-1}$ form an $F$-basis, interpreted
as the empty basis if $n=0$. Thus the quotient has dimension $n$;
if $|F|=q$, it has $q^n$ elements.
\end{theorem}
\begin{proof}
Divide a representative $g$ by $f$ to get $g=af+r$ with
$\deg r<n$. Thus $g$ and $r$ represent the same class.
If two such remainders differ by $bf$, their difference has
degree less than $n$. A nonzero $b$ would make the product have
degree at least $n$, so $b=0$ and the remainders are equal.
The scalar action on the quotient is multiplication by constant
classes, and the vector-space identities follow from ring
distributivity and the constants homomorphism. Each normal form
is a unique sum $\sum_{i=0}^{n-1}a_iX^i$, proving spanning
and independence of the asserted classes. Coordinate uniqueness
identifies the underlying set with $F^n$, proving both dimension
and the finite cardinality formula. For $n=0$ the only remainder
is zero, so the quotient is a one-element vector space and
$q^0=1$ is consistent.
\end{proof}
If $f=0$, no finite degree $n$ can be assigned by this theorem,
and $F[X]$ instead has the infinite basis $1,X,X^2,\ldots$.
An irreducible $f$ gives a field extension of degree $\deg f$;
irreducibility is unnecessary for the vector-space and counting
assertions.

\KEntry{K07.11.polynomial_function}{A polynomial and its function}
Let $F$ be a finite field with $q$ elements. For $a\ne0$,
multiplication by $a$ permutes the $q-1$ nonzero elements.
The product $P$ of these elements is nonzero, and comparison of
the products before and after the permutation gives
$P=a^{q-1}P$. Cancellation yields $a^{q-1}=1$.
It follows also at $a=0$ that $a^q-a=0$ for all $a\in F$.
However $X^q-X$ is a nonzero formal polynomial: $q\ge2$, and its
coefficient of index $q$ is $1$ while its coefficient of index one
is $-1$. Hence evaluation to the ring of functions $F\to F$ is
not injective. This proof uses only the field axioms and the finite
set of nonzero elements, and does not presuppose finite-field
classification.

Over an infinite domain, evaluation on all elements is injective.
If $f,g$ define the same function, $f-g$ vanishes everywhere.
A nonzero $f-g$ has at most $\deg(f-g)$ roots by
Theorem~\ref{thm:polynomial-root-bound}, contradicting infinitude.
Thus $f=g$. The distinction is between formal coefficient equality
and equality of a particular evaluation function; even where these
coincide they are different definitions.

\KEntry{K07.11.multivariate_boundary}{Several indeterminates}
A polynomial in $m$ commuting indeterminates over a commutative ring
has finite support on exponent vectors $\N^m$. Its monomials are
$X_1^{i_1}\cdots X_m^{i_m}$; coefficient addition and convolution
give the same ring-law proofs as in one variable.
For two variables the correspondence
\[
 \sum_{i,j}a_{ij}X^iY^j
 \longleftrightarrow
 \sum_j\left(\sum_i a_{ij}X^i\right)Y^j
\]
is a ring isomorphism $R[X,Y]\cong(R[X])[Y]$.
Finite support in one description is precisely finite support
in the other, and both product coefficients are
$\sum_{i+i'=k,\ j+j'=\ell}a_{ij}b_{i'j'}$.
This verifies the isomorphism, rather than treating the symbols
as an unexplained identification.
Total degree uses the sum of the exponent coordinates; degree
in $X$ or $Y$ uses one coordinate. Over a nonzero ring, for $X+Y^2$ these degrees
are respectively $2,1,2$. A lexicographic leading monomial with
$X>Y$ is $X$, and is not a total-degree assertion.

By Theorem~\ref{thm:polynomial-ufd}, $F[X,Y]$ is a UFD for a field
$F$, but it is not a PID. The ideal $(X,Y)$ is proper, because
evaluation at $(0,0)$ kills it but not $1$.
The elements $X,Y$ are prime: quotienting by $X$ gives the domain
$F[Y]$, and quotienting by $Y$ gives $F[X]$.
They are not associates, since $X$ does not divide $Y$, as is
seen by setting $X=0$. If $(X,Y)=(h)$, the nonzero $h$ divides
both $X$ and $Y$. Unique factorization and the primality of $X$
force such a common divisor to be a unit. This would make the
ideal the whole ring, a contradiction.

Division by an arbitrary list has no unrestricted analogue of the
unique univariate remainder. For example, order monomials
lexicographically with $X>Y$, and use the divisors
$g_1=XY-1$, $g_2=Y^2-1$. Starting with $XY^2$, cancellation
using $g_1$ gives
$XY^2-Yg_1=Y$; cancellation using $g_2$ gives
$XY^2-Xg_2=X$. Neither monomial $X$ nor $Y$ is divisible
by the leading monomials $XY$ or $Y^2$, so both are terminal
remainders under such a division rule, yet they differ.
Indeed $X-Y=Yg_1-Xg_2$ lies in the ideal.
Monomial orders, suitable generating sets, termination,
and uniqueness in multivariable algorithms belong to the
bounded transition toward D22; the one-variable division proof
does not silently provide those assertions.

\KEntry{K07.11.linear_algebra_bridge}{Operators and their polynomials}
Let $V$ be a vector space over a field $F$ and $T\in\End_F(V)$.
The map $a\mapsto a\id_V$ is a unital homomorphism to
$\End_F(V)$, with scalar maps commuting with $T$.
Theorem~\ref{thm:polynomial-evaluation} therefore defines $f(T)$,
with $(fg)(T)=f(T)g(T)$. Its kernel
\[
 I_T=\{f\in F[X]:f(T)=0\}
\]
is an ideal. If $I_T\ne0$, its unique monic generator is the
minimal polynomial $\mu_T$. Existence follows from the PID
property and scaling a nonzero generator to be monic; uniqueness
follows because generators differ by a unit, and two monic
associates are equal. The generator property proves
\[
 f(T)=0\quad\Longleftrightarrow\quad \mu_T\mid f.
\]
In particular it proves minimality of degree among nonzero
annihilating polynomials, since a nonzero multiple has at least
the generator's degree.

If $\dim_FV=d<\infty$, the D05 dimension theorem gives
$\dim_F\End_F(V)=d^2$, by the matrix units in a chosen basis.
Thus $\id,T,\ldots,T^{d^2}$ are linearly dependent, giving
a nonzero polynomial in $I_T$ and hence existence of $\mu_T$.
The characteristic polynomial is
$\chi_T(X)=\det(XI-[T])$, independent of basis because
$XI-P^{-1}[T]P=P^{-1}(XI-[T])P$ and determinants multiply.
We reuse the D05 Cayley--Hamilton theorem, with its
finite-dimensional hypothesis: $\chi_T(T)=0$.
The generator property just proved then gives
$\mu_T\mid\chi_T$ and $\deg\mu_T\le d$.
These are applications of the canonical D05 determinant and
Cayley--Hamilton results, not a second proof of them here.
For $V=0$, $\End_F(V)$ is the zero ring, $I_T=F[X]$,
and both polynomials are $1$ with degree zero.
For an infinite-dimensional space a minimal polynomial need
not exist: on the space with basis $e_0,e_1,\ldots$, the shift
$Te_i=e_{i+1}$ has $f(T)e_0=\sum a_ie_i\ne0$ for every
nonzero polynomial $f$.
Chapter~\ref{chap:M07.16} will express finite-dimensional operator
classification through the same $F[X]$-action, using only these
basic D05 facts and the PID structure theorem proved there.
